\subsection{Invertible formal power series}

PROPOSITION 5. --- In the ring A{[}{[}T{]}{]} of formal power series in one indeterminate the polynomial 1 - T is invertible, and we have \((1 - \mathrm{T})^{-1} = \sum_{n=0}^{\infty} \mathrm{T}^{n}\) .

For

\[
(1 - \mathrm{T}) \left(\sum_ {n = 0} ^ {\infty} \mathrm{T} ^ {n}\right) = \sum_ {n = 0} ^ {\infty} \mathrm{T} ^ {n} - \sum_ {n = 0} ^ {\infty} \mathrm{T} ^ {n + 1} = 1.
\]

PROPOSITION 6. --- Let \(u \in \mathbf{A}[[\mathbf{I}]]\) ; then for \(u\) to be invertible in \(\mathbf{A}[[\mathbf{T}]]\) it is necessary and sufficient that its constant term should be invertible in \(\mathbf{A}\) .

Suppose that there exists \(v \in A[[I]]\) such that \(uv = 1\) . Let \(\alpha, \beta\) be the constant terms of \(u\) and \(v\) , then \(\alpha\beta = 1\) , so \(\alpha\) is invertible.

Conversely, suppose that the constant term \(\alpha\) of \(u\) is invertible. Then there exists a formal power series \(t \in \mathbf{A}[[\mathrm{I}]]\) such that \(u = \alpha(1 - t)\) and \(\omega(t) > 0\) . Now there is a ring homomorphism \(\varphi: \mathbf{A}[[\mathrm{T}]] \to \mathbf{A}[[\mathrm{I}]]\) such that \(\varphi(\mathrm{T}) = t\) , and \(1 - \mathrm{T}\) is invertible in \(\mathbf{A}[[\mathrm{T}]]\) (Prop. 5); consequently \(1 - t\) is invertible in \(\mathbf{A}[[\mathrm{I}]]\) , and hence so is \(u\) .

Remark. --- Let M be the set of all formal power series with constant term 1. By Prop. 6, M is a commutative group under multiplication; the multiplicative group of A{[}{[}I{]}{]} is thus the direct product of M and the multiplicative group of A. We shall equip M with the topology induced from that of A{[}{[}I{]}{]}. For each \(\beta \in \mathbf{N}^{(I)}\) we have in IV, p.~26 defined the ideal \(a_{\beta}\) of A{[}{[}I{]}{]}; then \(1 + a_{\beta}\) is a subgroup of M and the family \((1 + a_{\beta})\) is a fundamental system of neighbourhoods of 1 in M. Since the multiplication in M is continuous, we see that M is a topological group (Gen.~Top., III, p.~223); in other words, the mapping \(f \mapsto f^{-1}\) is continuous in M.

Let K be a commutative field and \(\mathfrak{D}\) the subring of the field of rational fractions \(\mathrm{K}((\mathbf{X}_i)_{i\in I})\) formed of rational fractions in which the element 0 of \(\mathbf{K}^I\) is substitutable. If \(f\in \mathfrak{D}\) , we have \(f = \frac{u}{v}\) , where \(u\) and \(v\) are polynomials such that the constant term of \(v\) is \(\neq 0\) , hence \(v\) is invertible in \(\mathbf{K}[[\mathrm{I}]]\) . We can verify at once that the element \(uv^{-1}\) of \(\mathbf{K}[[\mathrm{I}]]\) depends only on \(f\) ; we say that the formal power series \(uv^{-1}\) is the expansion at the origin of the rational fraction \(\frac{u}{v}\) . The mapping \(f\mapsto uv^{-1}\) is an injective homomorphism of \(\mathfrak{D}\) into \(\mathbf{K}[[\mathrm{I}]]\) ; we shall often identify \(\mathfrak{D}\) with its image under this mapping.

